\documentclass[11pt,a4paper]{article}
\usepackage[margin=2.6cm]{geometry}
\usepackage[T1]{fontenc}
\usepackage{lmodern}
\setlength{\parindent}{0pt}
\setlength{\parskip}{0.9em}
\pagestyle{empty}

\begin{document}

14 September 2026

Editors\\
\textit{Complex \& Intelligent Systems}

Dear Editors,

Please consider our manuscript, ``Query Dual Dynamics Delta: Read-Driven Modulation of Delta-Rule Writes in Fixed-Capacity Fast-Weight Memory,'' as an Original Article in \textit{Complex \& Intelligent Systems}.

Fixed-capacity memories in adaptive intelligent systems must decide not only how to store an incoming association, but where limited correction effort will be useful when the state is later queried. We introduce Query Dual Dynamics Delta (QD$^3$), an online operator that summarizes historical queries in a compact low-rank state and uses that signal to modulate subsequent residual Delta-rule writes. The construction remains constant in sequence length, reduces exactly to ordinary Delta when modulation is disabled, and includes a clipped update with a simple local stability interpretation.

The experiments isolate this mechanism across controlled spatial data, text-derived and graph-derived workloads, and held-out WikiText-2 sequences. Validation-separated paired tests, matched write-trace and dual-trace controls, uniform-query conditions, and changing-demand experiments distinguish information supplied by read history from effects of added state or adaptive rates. Source code, prepared arrays, per-seed results, and run manifests accompany the submission.

We believe the work fits the journal's interest in machine learning, complex adaptive systems, and brain-like computing. It provides both a new memory operator and an interpretable account of how two interacting event streams---reads and writes---can guide resource allocation in a continuously updated intelligent system.

The manuscript is original, is not under consideration elsewhere, and has been approved by both authors. All figures and text are original to this submission.

Thank you for your consideration.

Sincerely,

Dian Yu\\
Corresponding author, on behalf of both authors

\end{document}
